\subsection{Prospective confirmation and merge abstention}
\label{sec:fresh-design}
We tested the follow-up's conditional findings on previously unused target histories, with no additional training or threshold selection. Before selecting these histories or computing repair outcomes, a separate protocol fixed all 36 full-information models, both calibration targets, the budget grid, hypotheses and analysis. The cohort construction above excludes every earlier target-history partition. No new organization, source schema or natural-error annotation is introduced.

A fixed second hash divides each source into eight disjoint blocks of 128 complete root histories, matching the exploratory evaluation size. All records of each selected root are retained, and each block is repaired separately. Inference context therefore remains 128 roots; full-log scaling is outside this experiment. Each block receives three new mixed-fault realizations, paired with the three previously fitted model seeds. Every budget and training arm sees the identical damaged graph within a block/seed. The complete experiment has 576 guarded budget/threshold conditions, plus 192 safety conditions at budget 64 and $q=.99$. No-edit outcomes are retained for all 48 block/seed inputs. No additional generated examples or verified fitting/calibration histories are supplied.

All confirmation policies use one simultaneous edit pass with atomic root substitutions. The safety comparison reuses identical candidates, scores and calibrated thresholds. Guarded decoding rejects tied additions and permits a merge only at score at least .99, without co-occurring handles, and when it is mutual best with a margin of at least .05 over alternatives. The \emph{never-merge} policy prohibits all alias merges while preserving the guarded policy's tied-addition rejection and atomic root substitutions. It can still repair participation and can still introduce wrong links. An unguarded policy retains the atomic-root contract but removes the additional confidence, margin and co-occurrence checks. Thus the comparison tests whether conservative merging adds value beyond complete merge abstention.

Two directional hypotheses were fixed. H1 predicts that augmentation at 64 bundles and $q=.99$ increases loan full-event recovery without increasing mean new damage. H2 predicts that guarded merging retains true identity pairs and improves history F1 or recovery over never merge without increasing mean new damage, tested separately for each source/arm. For uncertainty, we first average the three fixed model/fault runs within each block, then bootstrap eight paired blocks with 10,000 resamples. We report conditional 95\% percentile intervals. They describe new-history variation for these fixed models in these processes; they do not treat events or seeds as independent populations. Interval support requires a positive lower bound for the gain and a nonpositive upper bound for damage. Secondary comparisons have no multiplicity adjustment.

